\section{Introduction}

Dependently typed programming languages, such as Idris~\cite{idris-jfp},
Agda~\cite{norell2007thesis}, and Haskell with the appropriate extensions 
enabled~\cite{dephaskell}, allow us to give precise types which can
describe assumptions about and relationships between inputs and outputs.
This is valuable for reasoning about \emph{functional} properties, such as
correctness of algorithms on collections~\cite{McBride2014}, termination of
parsing~\cite{Danielsson2010} and 
scope safety of programs~\cite{allais-cpp17}.
However, reasoning about \emph{non-functional} properties in this setting,
such as memory safety, protocol correctness, or resource safety in general, is
more difficult though it can be achieved with techniques such as embedded
domain specific languages~\cite{brady-tfp14} or indexed
monads~\cite{Atkey2009,McBride2011}. These are, nevertheless, heavyweight
techniques which can be hard to compose.

Substructural type systems, such as linear type
systems~\cite{wadler-linear,Morris2016,linhaskell},
allow us to express \emph{when} an operation can be executed, by requiring
that a linear resource be accessed \emph{exactly once}. Being able to
combine linear and dependent types would give us the ability to express
an ordering on operations, as required by a protocol, with precision
on exactly what operations are allowed, at what time. Historically, however,
a difficulty in combining linear and dependent types has been in deciding
how to treat occurrences of variables in \emph{types}. This can be
avoided~\cite{Krishnaswami-ldtp} by never allowing types to depend on a
linear term, but more recent work on Quantitative Type Theory 
(QTT)~\cite{qtt,nuttin} solves the problem by assigning a \emph{quantity} to
each binder, and checking terms at a specific \emph{multiplicity}. 
Informally, in QTT, variables and function arguments have a multiplicity: $0$,
$1$ or unrestricted ($\omega$). We can freely use any variable in an argument
with multiplicity $0$---e.g., in types---but we can not use a variable with
multiplicity $0$ in an argument with any other
multiplicity. A variable bound with multiplicity $1$ must be used exactly once.
In this way, we can describe linear resource usage protocols, and furthermore
clearly express \emph{erasure} properties in types.
%We will make the concept
%of multiplicities more precise later in the paper.

Idris 2 is a new implementation of Idris, which uses QTT as its core
type theory. In this paper, we explore the possibilities of programming with
a full-scale language built on QTT. By full-scale, we mean a language with high
level features such as inference, interfaces, local function definitions
and other syntactic sugar. As an example, we will show how to implement a
library for concurrent programming with session types~\cite{Honda1993}.
We choose this example because, as demonstrated by the extensive literature
on the topic, correct concurrent programming is both hard to achieve, and
vital to the success of modern software engineering. Our aim is to show that
a language based on QTT is an ideal environment in
which to implement accessible tools for software developers, based on
decades of theoretical results.

%The code is
%submitted as anonymised supplementary material (\texttt{idris2-code.tgz}).

\subsection{Contributions}

This paper is about exploring what is possible in a language based on
Quantitative Type Theory (QTT), and introduces a new implementation of Idris.
We make the following research contributions:

\begin{itemize}
\item We describe Idris 2 (Section~\ref{sec:idris}), the first implementation
of quantitative type theory in a full programming language, and the first
language with full first-class dependent types implemented in itself.
\item We show how Idris 2 supports two important applications of quantities:
\emph{erasure} (Section~\ref{sec:erasure}) which gives type-level guarantees as
to which values are required at run-time, and \emph{linearity}
(Section~\ref{sec:linearity}) which gives type-level guarantees of resource
usage. We also describe a general approach to implementing linear
resource usage protocols (Section~\ref{sec:protocols}).
\item We give an example of QTT in practice, encoding bidirectional
session types (Section~\ref{sec:sessions}) for safe concurrent programming.
\end{itemize}

We do not discuss the metatheory of QTT, nor the trade-offs in its design
in any detail. Instead, our interest is in discussing how it has affected
the design of Idris 2, and in investigating the new kinds of programming and
reasoning it enables. Where appropriate, we will discuss the intuition behind
how argument multiplicities work in practice.

